\documentclass{article}
\usepackage{graphicx} % Required for inserting images
\usepackage{url}
\usepackage{listings}
\usepackage{listings-rust}
\usepackage{multirow}
\usepackage{xcolor}
\usepackage{float}
\usepackage[a4paper, left=2cm, right=2cm, top=3cm, bottom=3cm]{geometry}
\usepackage{array}
 
\definecolor{codegreen}{rgb}{0,0.6,0}
\definecolor{codegray}{rgb}{0.5,0.5,0.5}
\definecolor{codepurple}{rgb}{0.58,0,0.82}
\definecolor{backcolour}{rgb}{0.95,0.95,0.92}
 
\lstdefinestyle{mystyle}{
  backgroundcolor=\color{backcolour},
  commentstyle=\color{codegreen},
  keywordstyle=\color{magenta},
  numberstyle=\tiny\color{codegray},
  stringstyle=\color{codepurple},
  basicstyle=\ttfamily\footnotesize,
  breakatwhitespace=false,
  breaklines=true,
  captionpos=b,
  keepspaces=true,
  numbers=left,
  numbersep=5pt,
  showspaces=false,
  showstringspaces=false,
  showtabs=false,
  tabsize=2}
\lstset{style=mystyle}

\newcommand{\subsubsubsection}[1]{%
    \paragraph{#1} % Use paragraph for the subsubsubsection
    \mbox{} % Add a line break
}

\title{x86 Assembly}
\author{William Bland}
% \date{April 2025}

\begin{document}

\maketitle

\section{What is x86 assembly?}
Processors run by executing binary instructions called machine code, this machine code can be made into
a human readable format called assembly which can be turned back into machine code using a program called an assembler.
Each processor has an architecture, this architecture describes how the processor executes the machine code (so different architectures execute code differently,
so have different machine code and assembly languages).
For multiple processors to be able to execute the same machine code they must have the same architecture, some popular architectures include ARM, PowerPC, RISC-V and x86.
x86 is an example of a Complex Instruction Set Computer (CISC), this means that there are lots of instructions
so a lot cam be done in a single instruction so is very fast, this however uses a lot of power so is the main architecture in desktop computers.
This is in opposition to ARM which is a Reduced Instruction Set Computer (RISC) which has less instructions so has to execute many to
perform the same task so is slower than x86, this however uses less power so is the main architecture in mobile devices and embedded systems.

\section{Basics of x86 assembly}
x86 assembly has two syntaxes: Intel and AT\&T, all code will be in Intex syntax unless stated otherwise. \\
The code consists of an instruction followed by the registers it uses, or immediate values or a memory location.
Also a label can be placed to be able to refer to a specific line.

\section{Instructions}

\section{Directives}

\section{Sections}

\subsection{Data section}
db, equ, \$ etc

\section{Memory}

\subsection{Memory sizes and representaion}
Processors operate on digital signals which can either be high or low and can be interpreted in binary as a 1 or a 0, a single 1 or 0 is a bit.

\begin{table}[H]
  \centering
  \caption{Data sizes}
  \begin{tabular}{c|c|c|c}
    \textbf{Name} & \textbf{Unit} & \textbf{Bits} & \textbf{Bytes} \\
    \hline
    bit & b & 1 & 1/8 \\
    nibble & N/A & 4 & 1/2 \\
    byte & B & 8 & 1 \\
    word & N/A & 16 & 2 \\
    dword (double word) & N/A & 32 & 4 \\
    qword (quad word) & N/A & 64 & 8
  \end{tabular}
\end{table}

\noindent
Furthermore, prefixes can be added to the units to multiply the amount it stores
\begin{table}[H]
  \centering
  \caption{Data prefixes}
  \begin{tabular}{c|c|c|c}
    \textbf{Denary name} & \textbf{Denary size} & \textbf{Binary name} & \textbf{Binary size} \\
    \hline
    K & $10^3$ & Ki & $2^{10}$ \\
    M & $10^6$ & Mi & $2^{20}$ \\
    G & $10^9$ & Gi & $2^{30}$ \\
    T & $10^{12}$ & Ti & $2^{40}$ \\
    P & $10^{15}$ & Pi & $2^{50}$
  \end{tabular}
\end{table}

\noindent
Lastly, all memory addresses point to a unique byte in memory therefore a multibyte integer (eg 64-bit) could be stored in two different orderings (its endianess). In Little-Endian (LE) ordering the least significant byte occupies the smaller address space, while in Big-Endian (BE) ordering the most significant byte occupies the smaller address space. x86 and most other computers uses LE while data transmitted over networks uses BE.

\subsection{Registers}
Resisters store small values in small regions of memory accessible to the CPU, x86 has many registers that include
% https://wiki.osdev.org/CPU_Registers_x86-64

The registers on the same line overlap in memory (eg the 32-bit register is located in the lower half of the 64-bit register)

\begin{table}[H]
  \centering
  \caption{General purpose registers}
  \begin{tabular}{c|c|c|c|c|l}
    \textbf{64-bit} & \textbf{32-bit} & \textbf{16-bit} & \textbf{Higher 8-bit} & \textbf{Lower 8-bit} & \textbf{Description} \\
    \hline
    RAX & EAX & AX & AH & AL & Accumulator \\
    RBX & EBX & BX & BH & BL & Base \\
    RCX & ECX & CX & CH & CL & Counter \\
    RDX & EDX & DX & DH & DL & Data (often extends the A register) \\
    RSI & ESI & SI & N/A & SIL & Source index for string operations \\
    RDI & EDI & DI & N/A & DIL & Destination index for string operations \\
    RSP & ESP & SP & N/A & SPL & Stack pointer \\
    RBP & EBP & BP & N/A & BPL & Base pointer (for stack frames) \\
    R8 & R8D & R8W & N/A & R8B & General purpose register (long mode) \\
    R9 & R9D & R9W & N/A & R9B & General purpose register (long mode) \\
    R10 & R10D & R10W & N/A & R10B & General purpose register (long mode) \\
    R11 & R11D & R11W & N/A & R11B & General purpose register (long mode) \\
    R12 & R12D & R12W & N/A & R12B & General purpose register (long mode) \\
    R13 & R13D & R13W & N/A & R13B & General purpose register (long mode) \\
    R14 & R14D & R14W & N/A & R14B & General purpose register (long mode) \\
    R15 & R15D & R15W & N/A & R15B & General purpose register (long mode)
  \end{tabular}
\end{table}

\begin{table}[H]
  \centering
  \caption{Pointer registers}
  \begin{tabular}{c|c|c|l}
    \textbf{64-bit} & \textbf{32-bit} & \textbf{16-bit} & \textbf{Description} \\
    \hline
    RIP & EIP & IP & The pointer to the current instruction
  \end{tabular}
\end{table}

\begin{table}[H]
  \centering
  \caption{Segment registers (All are 16-bit)}
  \begin{tabular}{r|l}
    \textbf{Register} & \textbf{Description} \\
    \hline
    CS & Code segment \\
    DS & Data segment \\
    SS & Stack segment \\
    ES & Extra segment (for string operations) \\
    FS & General purpose segment register (Intel 80386) \\
    GS & General purpose segment register (Intel 80386)
  \end{tabular}
\end{table}

\subsection{Flags}

\subsection{Stack}

\subsection{Heap}

\subsection{Code}

\subsection{Data}

\subsection{Cache}

\subsection{Memory locations}

\subsection{Memory paging, protection and segmentation}
swap, page faults

\section{Operating modes}
\begin{table}[H]
  \centering
  %\caption{}
  \begin{tabular}{r|p{0.65\linewidth}}
    \textbf{Name} & \textbf{Description} \\
    \hline
    Real mode & Also called real address mode, all addresses always correspond to physical (real) memory addresses. Memory addresses are 20 bit allowing for 1MiB of memory, they are referenced by the value in a segment register left shifted by 4 then added to an offset for example [es:di]. There is no memory protection in real mode, and all x86 CPU's start in real mode when reset for compatibility \\
    Unreal mode & Also called big real mode, it is a variant of real mode, it allows the access to 4GiB of memory \\
    Protected mode & Also called protected virtual address mode, it allows the use of segmentation, virtual memory, paging and safe multi-tasking. To enter protected mode the software must setup a descriptor table and enable the Protection Enable (PE) bit in CR0 (Intel 80286). \\
    Virtual 8086 mode & Also called virtual read mode, it allows real mode processes to run in protected mode. It is a hardware virtualisation technique and uses the 20-bit segmentation that real mode uses but is subject to the protected mode's memory paging. \\
    Long mode & The 64-bit OS can access 64-bit registers and instructions. 64-bit programs run in a sub-mode called 64-bit mode and 32-bit and 16-bit protected mode programs run in another sub-mode called compatibility mode, however real mode and virtual 8086 mode programs cannot run natively in long mode. \\
    x86 virtualisation mode & Mode that uses hardware assisted virtualisation. \\
    AIS and 8080 emulation mode & Rarely used operating modes found on only a few old CPUs.
  \end{tabular}
\end{table}

\section{Branch prediction}
Spectre/meltdown

\section{x86 architecture}

\section{Intel vs AT\&T syntax}

\section{Interrupts}

\subsection{Syscalls (on Linux)}

In order for a program to communicate with the OS it has to hand over control of the processor back to the OS.
In 32-bit programs it uses the interrupt \texttt{int 80h}, in 64-bit programs it uses the instruction \texttt{syscall}.
The parameters for the syscall are placed in specific registers.


% https://www.man7.org/linux/man-pages/man2/syscall.2.html
\begin{table}[H]
  \centering
  \caption{Registers used for syscall parameters (not all args are used for all syscalls)}
  \begin{tabular}{c|c|l}
    \textbf{i386 Registers} & \textbf{x86-64 Registers} & \textbf{Useage} \\
    \hline
    eax & rax & Syscall number (different in 32-bit and 64-bit programs) \\
    ebx & rdi & Arg1 \\
    ecx & rsi & Arg2 \\
    edx & rdx & Arg3 \\
    esi & r10 & Arg4 \\
    edi & r8 & Arg5 \\
    ebp & r9 & Arg6
  \end{tabular}
\end{table}

% https://www.man7.org/linux/man-pages/man2/syscalls.2.html
% https://github.com/torvalds/linux/tree/master/arch/x86/entry/syscalls
\begin{table}[H]
  \centering
  \caption{Common Linux Syscalls}
  \begin{tabular}{|*{9}{>{\centering\arraybackslash}p{0.085\linewidth}|}}
    \hline \textbf{Name} & \textbf{eax} & \textbf{rax} & \textbf{Arg1} & \textbf{Arg2} & \textbf{Arg3} & \textbf{Arg4} & \textbf{Arg5} & \textbf{Arg6} \\
    \hline \texttt{write} & 4 & 1 & fd & ptr & count & - & - & - \\
    \hline \multicolumn{9}{|p{0.97\linewidth}|}{Writes the buffer (count bytes pointed to by ptr) to the file} \\
    \hline \texttt{exit} & 1 & 60 & error & - & - & - & - & - \\
    \hline \multicolumn{9}{|p{0.97\linewidth}|}{Exits the program with the specified error code} \\
    \hline \texttt{munmap} & 91 & 11 & addr & length & - & - & - & - \\
    \hline \multicolumn{9}{|p{0.97\linewidth}|}{Frees the length bytes at the address addr, must exactly match its corresponding mmap call} \\
    \hline \texttt{mremap} & ? & ? & ? & ? & - & - & - & - \\
    \hline \multicolumn{9}{|p{0.97\linewidth}|}{???} \\
    \hline \texttt{mmap} & 90 & 9 & addr & length & prot & flags & fd & offset \\
    \hline \multicolumn{9}{|p{0.97\linewidth}|}{Allocates length (must be a multiple of the page size [4096]) bytes, if addr is NULL (0) then any location is chosen else it will attempt to allocate at addr; prot is the protection flags for the memory, it is an bitwise OR of READ (1), WRITE (2), EXECUTE (4). Flags contains the mapping flags which is a bitwise OR of each flag such as MAP\_PRIVATE (0x2, changes are private to the calling process and are not written back) and MAP\_ANONYMOUS (0x20, allocates from main memory so not backed by a file as fd is ignored), fd can be set to -1 for anonymous mappings and offset is the offset from the start of the file (must be a multiple of the page size). eax/rax contains the ptr to the memory location allocated or an error between -4096 and -1. If a fd is given, the file is added to the programs virtual address space without being loaded into main memory (similar to swap memory). Additionally huge pages can be allocated with certain mappings and caution should be used when using this to access a file since race conditions can occur between processes if the file is not locked and the changes are not immediately written so use msync or fsync} \\
    \hline
  \end{tabular}
\end{table}

A fd (file descriptor) in Linux is an integer used to refer to a file, when opening a file the kernel returns an integer to the program to refer to the file with, and the kernel keeps track of the open file. Integers 0, 1 and 2 refer to stdin, stdout and stderr respectively, any negative fd is invalid so to specify no file, -1 is usually used.

\section{Appendix}

\subsection{32-bit Linux syscall}
\lstinputlisting{../linux_syscalls/32.asm}

\subsection{64-bit Linux syscall}
\lstinputlisting{../linux_syscalls/64.asm}

\section{Index}

\end{document}
